\include{zLaTex-Preamble}

\title{Homework 5} 
\author{Joshua Miller}
\date{}


%% \newcommand\invisiblesection[1]{%
%%   \refstepcounter{section}%
%%   \addcontentsline{toc}{section}{\protect\numberline{\thesection}#1}%
%%   \sectionmark{#1}}

\lhead{\fancyplain{}{{\Large\textbf{CMSC25400 / STAT 27725}} \\ 
    \large University of Chicago, Winter 2014  }}

\begin{document}
\rhead{\fancyplain{}{{\Large \textbf{Homework 5}}} \\ \large Joshua Miller}
% \maketitle
\vspace{40pt}~

\renewcommand{\m}[1]{\mathbf{#1}}

\section{Problem 1}

In order to test that the vectors produced by my code were actually
eigenvectors of $\m{A}$, I verified that for
$\epsilon' \approx \epsilon$ given as a parameter
\al{
\sum_i \lp (\m{A}v)_i - \lambda v_i \rp^2 < \epsilon'
}

The following bash commands are an overview of how to reproduce the results that follow. 
\lang{bash}
\begin{code}
# To build the programs, simply run make:

.../hw5$  make

# All three programs have usage statements. Examples are as follows

.../hw5$  ./eigP -e 1e-4 -k8 -f sample_eig.dat

# will find the 8 evecs with largest evals from input file sample.dat
# with error < 1e-4.  eigenvalues and eigenvectors are written to
# respective *.out files

.../hw5$  ./eigP -e 1e-4 -k8 -r 100

# the same as above but with a random matrix of size N = 100 //$
\end{code}

Figure~\ref{plot1} demonstrates the pattern of convergence given
varialbes $\epsilon$ and $N$ for matrix $\m A$. As one might expect,
the number of iterations required for the power method to converge is
proportional to the selected precision.  However it is also
interesting to note that the number of iterations to
convergence \emph{decreases} as the matrix size increases. 

\figH {plot1.png}{Plot of convergence iterations vs $\epsilon$ for
       selected matrix sizes}{plot1}{.8}

It can be
noted that the number of iterations grew logarithmically with
$\epsilon$, i.e. for the 10$\times$10 matrix,
\al{
\text{Iterations} \approx 0.780696\times\log(\epsilon)-0.130952
}

with an $r^2 = .9989$. Furthermore, as a function of matrix size for a
fized $\epsilon$, the number of iteration necessary as a function of matrix size $N$ was close to the following,
\al{
 \text{Iterations} \approx 49.5\times N^{-.43} + 9
}

Therefore, we can say that as a function of $\epsilon$, convergence is
on the order of $\mathcal{O}(\log(\epsilon))$.

\figH {plot2.png}{Plot of Frobenius norm vs $k$ precision
       level $\epsilon = 1\times 10^{-10}$ }{plot1}{.8}




\end{document}
